\section{Results}

\subsection{Ceilings}

The theoretical and measured ceilings for this card are collected in the table
below. The theoretical numbers are derived from device properties; the measured
numbers come from the FMA microbenchmark and the large device to device copy.

\inputgenerated{tables/ceilings.tex}

\subsection{The roofline}

Figure~\ref{fig:roofline-main} places every kernel configuration on the roofline
against both ceilings. Points whose intensity was computed from measured DRAM
bytes are drawn differently from points that fall back to a theoretical byte
count, so the reader always knows which is which.

% [H] pins the figure exactly here. Left as a floating [htbp] it was pushed to a
% page of its own and left the remainder of the page blank behind it.
\begin{figure}[H]
  \centering
  \includegenerated{roofline_main.pdf}
  \caption{Kernel ladder on the RTX 5070 roofline, theoretical and measured
    ceilings shown together. One label per kernel family gives the largest
    problem size that family was run at.}
  \label{fig:roofline-main}
\end{figure}

\subsection{Nsight Compute counters}

The counters that explain each kernel's position, achieved occupancy, DRAM
throughput, shared memory bank conflicts, and coalescing measured as sectors per
request, are collected for the curated subset below. A perfectly coalesced warp
moves 4 sectors per request for a 4 byte access and 16 for a \texttt{float4}, so
the ideal depends on the access width and is not a single number.

\inputgeneratedwide{tables/ncu_counters.tex}

\subsection{Timing summary}

Every configuration in the sweep. The table runs over several pages.

\inputgenerated{tables/timing_summary.tex}
